\section{Resultados}\label{sec:resultados}

\subsection{Caso de la profesora}

Caso \texttt{casos/entrada\_G2\_red\_15\_nodos.json}, tal como se recibió: $N=15$, $k=5$, $B=30$ y
$u_i(x)=w_i\ln(1+x)$, con $w_i$ el coeficiente de cada nodo dado por la profesora. La entrada trae
solo los coeficientes, así que las categorías de datos son \textbf{sintéticas}: el programa las
genera al leer el caso con semilla 36921, universo de 20 categorías y entre 3 y 7 por nodo, sin
modificar el archivo. La salida lo declara en \texttt{origen\_categorias}, y
\texttt{src/completar\_categorias.py} escribe esas mismas categorías en un archivo si se quieren ver.

\begin{center}
\begin{tabular}{>{\raggedright}p{0.36\linewidth}>{\raggedright\arraybackslash}p{0.52\linewidth}}
\cabecera \cab{Resultado} & \cab{Valor} \\
Nodos elegidos (en orden) & N04, N15, N14, N13, N06 \\
Ganancias marginales $\Delta$ de la selección & 7, 6, 4, 2, 1 \\
Categorías cubiertas $f(A)$ & 20 de 20 \\
Óptimo exacto & 20, con $\binom{15}{5}=3003$ evaluaciones; razón 1.0000 \\
Reparto de $B=30$ & N04 $=7$, N15 $=5$, N14 $=6$, N13 $=5$, N06 $=7$ \\
Utilidad total & 99.5821, igual al óptimo por programación dinámica para esos cinco nodos \\
Tiempo (selección $+$ reparto) & 0.018 ms (mediana de 5; depende del equipo) \\
Verificación & 12 de 12 chequeos \\
\bottomrule
\end{tabular}
\end{center}

\paragraph{Traza de la selección} (\texttt{resultados/entrada\_G2\_red\_15\_nodos/traza.md}):

\begin{center}
\begin{tabular}{ccc>{\raggedright}p{0.42\linewidth}c}
\cabecera \cab{Paso} & \cab{Elegido} & \cab{$\Delta$} & \cab{Categorías nuevas} & \cab{$f(A)$} \\
1 & N04 & 7 & K01, K02, K06, K09, K12, K13, K14 & 7 \\
2 & N15 & 6 & K03, K08, K10, K16, K17, K19 & 13 \\
3 & N14 & 4 & K04, K07, K11, K15 & 17 \\
4 & N13 & 2 & K05, K20 & 19 \\
5 & N06 & 1 & K18 & 20 \\
\bottomrule
\end{tabular}
\end{center}

Los empates se resuelven por el coeficiente: en el paso 1 empatan con 7 N04, N06, N13, N14 y N15,
y gana N04 ($w=12.0$); en el paso 2 empatan N13 y N15 con 6, y gana N15 ($w=9.5$ frente a 8.3); en
el paso 5 empatan N05, N06 y N11 con 1, y gana N06 ($w=11.1$). Las ganancias nunca suben, como
predice la submodularidad.

\paragraph{Efecto del desempate.} 33 grupos de cinco nodos cubren las 20 categorías; tras el
reparto, su utilidad va de 91.06 a 107.04. La versión 1.0.0, que desempataba por posición en el
archivo, entregaba N04, N13, N14, N03 y N05 con utilidad 92.58; desempatando por coeficiente, la
cobertura es la misma y la utilidad sube a 99.58. La mejor de las 33 (107.04) no se alcanza porque
selección y reparto se resuelven en secuencia (Sección~\ref{sec:formalizacion}).

\paragraph{Traza del reparto.} La primera unidad va a N04 ($12.0\ln2=8.3178$), la segunda a N06
($11.1\ln2=7.6939$), luego N14 (7.0008), N15 (6.5849) y N13 (5.7531). La sexta vuelve a N04
($12.0(\ln3-\ln2)=4.8656$). Las 30 ganancias marginales aparecen en orden no creciente en
\texttt{traza.md} y \texttt{traza\_reparto.csv}.

\subsection{Casos sintéticos de ejemplo}

Casos de \texttt{casos/ejemplos/}, generados con semilla fija (\texttt{src/generar\_casos\_ejemplo.py}):

\begin{center}
\begin{tabular}{lrrrrlrr}
\cabecera \cab{Caso} & \cab{$N$} & \cab{$k$} & \cab{$B$} & \cab{$f$} & \cab{Razón vs. óptimo} & \cab{Utilidad} & \cab{ms} \\
ej00\_tres\_nodos & 3 & 2 & 6 & 5 & 0.8333 (5/6) & 14.0511 & 0.006 \\
ej01\_pequeno & 8 & 3 & 20 & 11 & 1.0000 & 55.3778 & 0.009 \\
ej02\_pequeno\_caps & 10 & 4 & 25 & 18 & 0.9474 (18/19) & 34.1613 & 0.012 \\
ej03\_mediano & 14 & 5 & 40 & 39 & 1.0000 & 80.4221 & 0.024 \\
ej04\_grande & 500 & 20 & 2000 & 400 & no se calcula & 1221.8066 & 1.750 \\
ej05\_grande\_caps & 2000 & 40 & 10000 & 800 & no se calcula & 955.9874 & 11.686 \\
\bottomrule
\end{tabular}
\end{center}

Todos cumplen los 12 chequeos. En ej04 y ej05 el voraz cubre todo el universo (400 y 800
categorías), así que ahí también es óptimo. En los casos con óptimo calculado, la menor razón es
0.8333, muy por encima de la garantía de 0.632. Con el desempate por la ganancia de la primera unidad, la cobertura
es la misma que con la versión 1.0.0 en todos los ejemplos y la utilidad no baja en ninguno (sube,
por ejemplo, de 1054.64 a 1221.81 en ej04).
